%%=============================================================================
%% Voorwoord
%%=============================================================================

\chapter*{\IfLanguageName{dutch}{Woord vooraf}{Preface}}%
\label{ch:voorwoord}

Dit onderzoek is voortgekomen uit een praktische frustatie: bij het opzetten van nieuwe frontend-projecten bij digital agencies zoals Ballistix Digital blijkt veel tijd verloren te gaan aan repetitieve taken. Hoewel Large Language Models nu steeds geavanceerder worden, bleek het helemaal niet vanzelfsprekend om ze in te zetten voor betrouwbare codegeneratie. Ik heb ervaren dat pure AI-ondersteuning al snel hallucinaties produceert en architecturale afspraken mist, wat het uiteindelijke voordeel volledig in het water doet vallen.

De vraag waar ik dus tegenaan liep was: hoe zet je AI op een manier in die werkelijk waarde toevoegt, zonder dat je halverwege toch terug bent bij afvinken en herstellen? Dit onderzoek probeert daar een antwoord op te geven door twee werkelijk verschillende benaderingen te vergelijken --- van volledig vrij redeneren tot een strakker ingerichte template-gebaseerde aanpak.

Dit werkstuk is niet alleen het resultaat van onderzoek, maar vooral van veel experimenteren, falen, en opnieuw beginnen --- precies wat het werken met AI lijkt te zijn. Ik hoop dat deze bachelorproef nut heeft voor wie worstelt met dezelfde vraag.

\vspace{\baselineskip}

Ouahab Yanis, mei 2026